\documentclass[11pt,a4paper]{article}
\usepackage{geometry}
\geometry{top=2.5cm,bottom=2.5cm,left=2.3cm,right=2.3cm}

\usepackage{xeCJK}
\setCJKmainfont{PingFang SC}[BoldFont=PingFang SC Semibold]
\setCJKsansfont{PingFang SC}
\setCJKmonofont{PingFang SC}

\setmainfont{Helvetica Neue}[BoldFont=Helvetica Neue Bold]
\setsansfont{Helvetica Neue}
\setmonofont{Menlo}[Scale=0.88]

\usepackage{xcolor}
\definecolor{navyblue}{HTML}{0F172A}
\definecolor{accentblue}{HTML}{0284C7}
\definecolor{borderblue}{HTML}{38BDF8}
\definecolor{bgblue}{HTML}{F0F9FF}

\definecolor{warnamber}{HTML}{D97706}
\definecolor{borderamber}{HTML}{F59E0B}
\definecolor{bgamber}{HTML}{FFFBEB}

\definecolor{dangerred}{HTML}{DC2626}
\definecolor{borderred}{HTML}{EF4444}
\definecolor{bgred}{HTML}{FEF2F2}

\definecolor{tipgreen}{HTML}{16A34A}
\definecolor{bordergreen}{HTML}{22C55E}
\definecolor{bggreen}{HTML}{F0FDF4}

\definecolor{noteslate}{HTML}{475569}
\definecolor{borderslate}{HTML}{94A3B8}
\definecolor{bgslate}{HTML}{F8FAFC}

\definecolor{codebg}{HTML}{F1F5F9}
\definecolor{codeframe}{HTML}{CBD5E1}

\usepackage{tcolorbox}
\tcbuselibrary{skins,breakable}

% Custom GFM Alert Boxes
\newtcolorbox{gfmImportant}[1][]{
  colback=bgblue,
  colframe=accentblue,
  leftrule=4pt,
  rightrule=0.5pt,
  toprule=0.5pt,
  bottomrule=0.5pt,
  arc=3pt,
  boxsep=4pt,
  left=8pt,right=8pt,top=6pt,bottom=6pt,
  breakable,
  title={\textbf{\color{navyblue} IMPORTANT #1}},
  coltitle=navyblue
}

\newtcolorbox{gfmWarning}[1][]{
  colback=bgamber,
  colframe=warnamber,
  leftrule=4pt,
  rightrule=0.5pt,
  toprule=0.5pt,
  bottomrule=0.5pt,
  arc=3pt,
  boxsep=4pt,
  left=8pt,right=8pt,top=6pt,bottom=6pt,
  breakable,
  title={\textbf{\color{warnamber} WARNING #1}},
  coltitle=warnamber
}

\newtcolorbox{gfmNote}[1][]{
  colback=bgslate,
  colframe=noteslate,
  leftrule=4pt,
  rightrule=0.5pt,
  toprule=0.5pt,
  bottomrule=0.5pt,
  arc=3pt,
  boxsep=4pt,
  left=8pt,right=8pt,top=6pt,bottom=6pt,
  breakable,
  title={\textbf{\color{noteslate} NOTE #1}},
  coltitle=noteslate
}

\newtcolorbox{gfmTip}[1][]{
  colback=bggreen,
  colframe=tipgreen,
  leftrule=4pt,
  rightrule=0.5pt,
  toprule=0.5pt,
  bottomrule=0.5pt,
  arc=3pt,
  boxsep=4pt,
  left=8pt,right=8pt,top=6pt,bottom=6pt,
  breakable,
  title={\textbf{\color{tipgreen} TIP #1}},
  coltitle=tipgreen
}

\newtcolorbox{gfmCaution}[1][]{
  colback=bgred,
  colframe=dangerred,
  leftrule=4pt,
  rightrule=0.5pt,
  toprule=0.5pt,
  bottomrule=0.5pt,
  arc=3pt,
  boxsep=4pt,
  left=8pt,right=8pt,top=6pt,bottom=6pt,
  breakable,
  title={\textbf{\color{dangerred} CAUTION #1}},
  coltitle=dangerred
}

\newtcolorbox{gfmBlockquote}[1][]{
  colback=bgslate,
  colframe=borderslate,
  leftrule=3pt,
  rightrule=0.2pt,
  toprule=0.2pt,
  bottomrule=0.2pt,
  arc=2pt,
  boxsep=3pt,
  left=6pt,right=6pt,top=4pt,bottom=4pt,
  breakable
}

\usepackage{tabularx}
\usepackage{booktabs}
\usepackage{colortbl}
\usepackage{multirow}

\usepackage{listings}
\lstset{
  basicstyle=\ttfamily\small,
  backgroundcolor=\color{codebg},
  frame=single,
  rulecolor=\color{codeframe},
  breaklines=true,
  breakatwhitespace=true,
  showstringspaces=false,
  tabsize=2,
  captionpos=b,
  keywordstyle=\color{accentblue}\bfseries,
  commentstyle=\color{noteslate}\itshape,
  stringstyle=\color{tipgreen}
}

\usepackage{graphicx}
\usepackage{fancyhdr}
\usepackage{lastpage}
\usepackage{hyperref}
\hypersetup{
  colorlinks=true,
  linkcolor=accentblue,
  urlcolor=accentblue,
  citecolor=accentblue,
  pdfborder={0 0 0}
}

\pagestyle{fancy}
\fancyhf{}
\fancyhead[L]{\small\color{noteslate}Atlas Architecture Design}
\fancyhead[R]{\small\color{noteslate}Atlas Architecture Specification}
\fancyfoot[C]{\small\color{noteslate}Page \thepage\ of \pageref{LastPage}}
\renewcommand{\headrulewidth}{0.4pt}
\renewcommand{\footrulewidth}{0.4pt}

\title{\textbf{\Large Atlas Architecture Design}}
\author{Atlas Architecture Group}
\date{\today}

\begin{document}
\maketitle
\tableofcontents
\newpage

\textbf{Version:} v1.0.2

\textbf{Status:} Frozen Production Baseline (Topology Conformance Amendment)

\textbf{Owner:} snkio027

\textbf{Scope:} Kubernetes Runtime, Argo CD GitOps Control Plane, Developer Platform Foundation, Streaming/Lakehouse Infrastructure

\vspace{0.5em}\hrule\vspace{0.5em}
\subsection{0. 规范性引用与文档治理}
Atlas 采用分层架构文档体系。不同层级的设计文档分别拥有不同的规范权威，以避免平台级原则、领域级实现语义与具体工程实现相互覆盖。

\subsubsection{0.1 Normative Architecture References}
\begin{enumerate}
  \item \textbf{Atlas Architecture Design v1.0.2（本文档）}
定义 Atlas 的平台级架构原则、信任边界、控制权模型、系统不变量、容灾边界与领域分层。

\end{enumerate}
\begin{enumerate}
  \item \textbf{Atlas GitOps Control Plane Design v1.0.3}
定义 GitOps 控制域内的具体执行语义，包括 External Root Anchor、AppProject 自举链、Two-Level Reconciliation DAG、Sync Wave、Application Health Gate、删除保护与 Break-Glass Recovery。

\end{enumerate}
\begin{enumerate}
  \item \textbf{AI-Native Platform Control Plane Standard v1.4}
定义 AI Agent 的交互权限、Human Judgment Gate、鉴权、审计与安全兜底模型。

\end{enumerate}
\subsubsection{0.2 规范优先级}
Atlas 明确区分：

\begin{itemize}
  \item \textbf{Global Architecture Invariant}：由本文档定义，领域设计不得违反。
  \item \textbf{Domain Execution Semantics}：由对应领域设计定义，在本文档授权范围内拥有更高的执行细节权威。
  \item \textbf{Implementation Detail}：由代码、Manifest、Helm Values、Policy 等具体工程资产实现。
\end{itemize}
因此：

\begin{gfmBlockquote}
Domain Design MAY refine Architecture semantics, but MUST NOT override Architecture invariants.
\end{gfmBlockquote}
当领域文档与本文档中的非规范性示例冲突时，以领域文档为准；当领域文档与本文档定义的平台级不变量冲突时，以本文档为准。

本文中的 \textbf{MUST / MUST NOT / SHOULD / SHOULD NOT / MAY} 具有规范性约束意义。

\vspace{0.5em}\hrule\vspace{0.5em}
\subsection{1. 架构愿景与定位}
Atlas 是面向云原生数据基础设施与流处理场景的内部开发者平台（Internal Developer Platform, IDP）参考架构。

其首要目标是在：

\begin{itemize}
  \item macOS Apple Silicon
  \item OrbStack Docker Runtime
  \item Kubernetes
  \item Argo CD
  \item Streaming / Lakehouse Infrastructure
\end{itemize}
组成的本地开发环境中，提供：

\begin{itemize}
  \item 可完全离线恢复的基础设施供应链；
  \item 极速、确定性的环境拉起能力；
  \item 严格声明式的稳态治理模型；
  \item 可预测的控制面恢复能力；
  \item 面向开发者的低认知负荷抽象；
  \item 面向 AI Agent 的显式权限边界与安全操作域。
\end{itemize}
Atlas 不将自身定位为若干云原生工具的集合，而将平台视为一个由多个独立 Reconciler 共同组成的\textbf{分层控制系统}。

平台设计遵循 CNCF Platform Engineering 的核心思想：

\begin{gfmBlockquote}
Platform as a Product
Self-Service
Declarative Automation
Cognitive Load Reduction
Clear Ownership
Failure Domain Isolation
\end{gfmBlockquote}
同时，Atlas 将 AI Agent 视为未来平台工程的一等参与者，因此物理目录、控制权、信任级别和执行权限必须能够被机器明确判断，而不能依赖隐含的人类经验。

\vspace{0.5em}\hrule\vspace{0.5em}
\subsection{2. 核心设计原则}
Atlas 建立在以下不可妥协的架构原则之上。

\vspace{0.5em}\hrule\vspace{0.5em}
\subsubsection{2.1 Bootstrap ≠ GitOps}
\_\_INLINE\_CODE\_0\_\_ 与 \_\_INLINE\_CODE\_1\_\_ 属于完全不同的生命周期域。

\paragraph{Bootstrap}
Bootstrap 是一次性、命令式、可退出的控制面建立过程。

其职责仅包括：

\begin{itemize}
  \item 宿主环境检查；
  \item Kubernetes substrate 建立；
  \item Local Artifact Registry 建立与预热；
  \item 极简 Argo CD Seed 安装；
  \item Tier-0 Bootstrap AppProject 建立；
  \item External Root Anchor 实例化。
\end{itemize}
Bootstrap MUST NOT 长期参与平台稳态资源管理。

完成 GitOps 控制权移交后：

\begin{gfmBlockquote}
Bootstrap MUST terminate as an active control authority.
\end{gfmBlockquote}
\paragraph{GitOps}
\_\_INLINE\_CODE\_0\_\_ 保存平台长期 Desired State。

Argo CD 负责持续比较：

\begin{lstlisting}[language=text]
Git Desired State
        ↓
Kubernetes Actual State
\end{lstlisting}
并持续进行收敛。

Bootstrap 与 GitOps 的关系不是双向共同管理，而是：

\begin{lstlisting}[language=text]
Bootstrap
   │
   │ instantiate
   ▼
GitOps Control Plane
   │
   │ adopt
   ▼
Steady State
\end{lstlisting}
控制权转移 MUST 是单向的。

\vspace{0.5em}\hrule\vspace{0.5em}
\subsubsection{2.2 四阶控制权模型}
Atlas 将系统控制权划分为四个不同的 Authority Domain。

\paragraph{Git owns Definition}
Git 是平台期望状态的定义权威。

Git 保存：

\begin{itemize}
  \item Root Anchor 契约模板；
  \item Platform Desired State；
  \item Workload Desired State；
  \item Policy；
  \item Environment Overlay；
  \item Version Lock；
  \item Governance Intent。
\end{itemize}
Git 回答：

\begin{gfmBlockquote}
系统“应该是什么”。
\end{gfmBlockquote}
\vspace{0.5em}\hrule\vspace{0.5em}
\paragraph{Bootstrap owns Instantiation}
Bootstrap 拥有初始实例化权。

它负责：

\begin{itemize}
  \item 读取环境上下文；
  \item 建立最小 Kubernetes / Argo Seed；
  \item 建立 Tier-0 Bootstrap Capability；
  \item 实例化 External Root Anchor。
\end{itemize}
Bootstrap 回答：

\begin{gfmBlockquote}
当 GitOps 本身尚不存在时，由谁启动 GitOps。
\end{gfmBlockquote}
\vspace{0.5em}\hrule\vspace{0.5em}
\paragraph{Argo CD owns GitOps Reconciliation}
Argo CD 是 Atlas GitOps Desired-State Domain 中唯一的持续同步权威。

它负责：

\begin{itemize}
  \item Git → Kubernetes 收敛；
  \item Application 拓扑；
  \item Sync；
  \item Prune；
  \item Health Gate；
  \item Drift Correction。
\end{itemize}
Argo CD 不拥有应用内部运行时行为。

\vspace{0.5em}\hrule\vspace{0.5em}
\paragraph{Operators own Runtime Behavior}
资源进入 Kubernetes API 后，其运行时生命周期由 Kubernetes 原生 Controller 或 Domain Operator 负责。

例如：

\begin{lstlisting}[language=text]
Git
 ↓
Argo CD
 ↓
FlinkDeployment CR
 ↓
Flink Kubernetes Operator
 ↓
JobManager / TaskManager
\end{lstlisting}
或者：

\begin{lstlisting}[language=text]
Git
 ↓
Argo CD
 ↓
Redpanda Cluster CR
 ↓
Redpanda Operator
 ↓
Runtime Cluster
\end{lstlisting}
因此：

\begin{gfmBlockquote}
GitOps Reconciliation Authority ≠ Runtime Reconciliation Authority.
\end{gfmBlockquote}
\_\_INLINE\_CODE\_0\_\_ 只是 Argo CD 对自身 Desired State 的自管理机制，不构成独立的第五控制域。

\vspace{0.5em}\hrule\vspace{0.5em}
\subsubsection{2.3 External Root Anchor + Two-Level Reconciliation DAG}
Atlas 正式采用：

\begin{gfmBlockquote}
\textbf{External Root Anchor + Two-Level Reconciliation DAG Pattern}
\end{gfmBlockquote}
作为 GitOps 拓扑的架构基线。

这是 Atlas v1.0.2 的核心拓扑不变量。

\vspace{0.5em}\hrule\vspace{0.5em}
\paragraph{2.3.1 External Root Anchor}
External Root Anchor 是 Atlas GitOps 控制图的外部信任锚点。

它：

\begin{itemize}
  \item 定义于 Git；
  \item 由 Bootstrap 根据环境上下文实例化；
  \item 作为 Root App-of-Apps 的入口；
  \item 不由另一个 Argo CD Application 持续调和。
\end{itemize}
其目的是显式终止无限自引用：

\begin{lstlisting}[language=text]
Who manages Argo CD?
        ↓
argocd-self

Who manages argocd-self?
        ↓
Platform Control DAG

Who manages Platform Control?
        ↓
External Root Anchor

Who manages External Root Anchor?
        ↓
Git Definition
+
Bootstrap Instantiation
+
Human / Policy Governance
\end{lstlisting}
External Root Anchor MUST NOT 被设计为自我管理 Application。

\vspace{0.5em}\hrule\vspace{0.5em}
\paragraph{2.3.2 Level 1 — Root Macro DAG}
Root 层只表达：

\begin{gfmBlockquote}
\textbf{Trust Transition 与 Macro Readiness Gate}
\end{gfmBlockquote}
Root MUST NOT 成为 Platform Component Catalog。

Root 不应知道：

\begin{itemize}
  \item Flink；
  \item Envoy Gateway；
  \item Redpanda；
  \item MinIO；
  \item Loki；
  \item Tempo；
  \item Alloy；
  \item 未来新增的普通平台组件。
\end{itemize}
Root 只负责少量稳定的控制域跃迁：

\begin{lstlisting}[language=text]
External Root Anchor
        │
        ├── Project Bootstrap
        │
        ├── Platform Control
        │
        └── Workload Control
\end{lstlisting}
因此普通平台能力增加或删除 SHOULD NOT 修改 \_\_INLINE\_CODE\_0\_\_。

\vspace{0.5em}\hrule\vspace{0.5em}
\paragraph{2.3.3 Level 2 — Platform Capability DAG}
平台组件的实际依赖关系由：

\begin{lstlisting}[language=text]
gitops/platform/applications/
\end{lstlisting}
表达。

Platform Control 负责：

\begin{itemize}
  \item Foundation；
  \item Management；
  \item Operators；
  \item Infrastructure Controllers；
  \item Platform Services。
\end{itemize}
典型逻辑依赖为：

\begin{lstlisting}[language=text]
Foundation
    ↓
Management
    ↓
Operators
    ↓
Controllers
    ↓
Platform Services
\end{lstlisting}
Root 仅关心：

\begin{lstlisting}[language=text]
Platform Control Healthy?
\end{lstlisting}
而不关心其内部拥有多少组件。

\vspace{0.5em}\hrule\vspace{0.5em}
\paragraph{2.3.4 Root manages trust transitions}
正式约束：

\begin{gfmBlockquote}
\textbf{Root SHALL model trust-domain transitions, not platform components.}
\end{gfmBlockquote}
任何只是因为新增：

\begin{itemize}
  \item 数据库；
  \item Operator；
  \item Observability Backend；
  \item Gateway Controller；
  \item Messaging Engine；
\end{itemize}
而要求修改 Tier-0 Root 的设计，默认视为架构异味。

\vspace{0.5em}\hrule\vspace{0.5em}
\paragraph{2.3.5 Leaf Application 是最小失败隔离单元}
Atlas 将独立 Argo CD Application 视为最小的：

\begin{itemize}
  \item Sync Unit；
  \item Health Unit；
  \item Autosync Unit；
  \item Recovery Unit；
  \item Failure Isolation Unit。
\end{itemize}
例如：

\begin{lstlisting}[language=text]
argocd-self
sealed-secrets
flink-operator
envoy-gateway
kube-prometheus-stack
redpanda
minio
\end{lstlisting}
SHOULD 分别成为独立 Leaf Application。

Application 不应仅因为文件系统存在子目录而进一步嵌套。

\vspace{0.5em}\hrule\vspace{0.5em}
\paragraph{2.3.6 Application 嵌套深度限制}
External Root Anchor 以下默认最多允许两个 Application orchestration hops：

\begin{lstlisting}[language=text]
External Root Anchor
        ↓
Control Application
        ↓
Leaf Application
        ↓
Kubernetes Resources / CRs
\end{lstlisting}
禁止形成：

\begin{lstlisting}[language=text]
Root
 ↓
Platform
 ↓
Domain
 ↓
Subdomain
 ↓
Component
 ↓
Resources
\end{lstlisting}
仅当存在真实的：

\begin{itemize}
  \item Trust Boundary；
  \item Ownership Boundary；
  \item Failure Isolation Boundary；
\end{itemize}
时才允许增加新的 Application 层级。

突破默认深度 MUST 通过 ADR 明确批准。

\vspace{0.5em}\hrule\vspace{0.5em}
\subsubsection{2.4 Directory Ownership ≠ Dependency Ordering}
Atlas 强制区分两个正交维度：

\paragraph{Directory}
回答：

\begin{gfmBlockquote}
这个能力属于哪个领域？
\end{gfmBlockquote}
\paragraph{Sync Wave}
回答：

\begin{gfmBlockquote}
这个能力在什么依赖条件满足后才允许收敛？
\end{gfmBlockquote}
因此：

\begin{lstlisting}[language=text]
management/
operators/
networking/
observability/
messaging/
storage/
\end{lstlisting}
表达领域归属。

而：

\begin{lstlisting}[language=text]
Management
Operators
Controllers
Platform Services
\end{lstlisting}
表达依赖阶段。

Atlas MUST NOT 创建：

\begin{lstlisting}[language=text]
controllers/
platform-services/
\end{lstlisting}
仅用于机械对应 Sync Wave。

同一个领域内允许存在不同依赖阶段。

例如：

\begin{lstlisting}[language=text]
networking/
├── envoy-gateway-controller
└── gateway-runtime-resources
\end{lstlisting}
可以分别处于不同 Wave。

\vspace{0.5em}\hrule\vspace{0.5em}
\subsubsection{2.5 GitOps Reconciliation Authority}
Argo CD 是 Atlas GitOps Desired-State Domain 的唯一持续同步权威。

平台禁止同时引入第二个与 Argo CD 对相同资源进行 Desired-State Reconciliation 的控制器。

例如禁止：

\begin{lstlisting}[language=text]
Argo CD
+
Helm Release State
+
另一个 GitOps Controller
\end{lstlisting}
共同竞争同一资源所有权。

Helm 在 Atlas 中仅作为：

\begin{gfmBlockquote}
Template Renderer
\end{gfmBlockquote}
使用。

Atlas MUST NOT 依赖 Helm Release History 作为平台稳态的一部分。

\vspace{0.5em}\hrule\vspace{0.5em}
\subsubsection{2.6 Artifacts ≠ Intent}
Atlas 将软件发行产物与平台治理意图彻底分离。

\paragraph{Artifacts}
第三方发行制品属于 Supply Chain Artifact。

例如：

\begin{itemize}
  \item Helm Chart \_\_INLINE\_CODE\_0\_\_；
  \item OCI Image；
  \item CRD Bundle；
  \item 其他不可变发行资产。
\end{itemize}
复杂第三方 Helm Chart SHOULD 以不可变形式缓存，并在 Root DAG 启动前完成本地可用性准备。

\paragraph{Intent}
Atlas 自身的治理意图属于 Git Desired State，例如：

\begin{lstlisting}[language=text]
gitops/**/values.yaml
resources/
overlays/
policies/
\end{lstlisting}
Argo CD 在渲染阶段将：

\begin{lstlisting}[language=text]
Immutable Upstream Artifact
          +
Atlas Governance Intent
          ↓
Hydrated Desired State
\end{lstlisting}
进行组合。

\vspace{0.5em}\hrule\vspace{0.5em}
\paragraph{2.6.1 Offline Supply Chain Invariant}
“完全离线”不只意味着 Helm Chart 本地可用。

Atlas 的离线不变量是：

\begin{gfmBlockquote}
Runtime 所需的全部外部 Artifact MUST 在断开公网前被本地验证并可用。
\end{gfmBlockquote}
因此离线供应链最终必须同时覆盖：

\begin{lstlisting}[language=text]
Helm / OCI Artifacts
+
Container Images
+
Version Metadata
+
Required CRDs
\end{lstlisting}
具体缓存、镜像同步和 Registry Distribution 机制由独立 Supply Chain Design 或 RFC 管辖。

\vspace{0.5em}\hrule\vspace{0.5em}
\subsubsection{2.7 Developer Cognitive Load Reduction}
Atlas 根据使用者角色分级暴露复杂性。

\paragraph{Platform Infrastructure}
平台基础设施允许使用 Helm。

原因是：

\begin{itemize}
  \item 上游生态复杂；
  \item 平台团队拥有治理复杂性的能力；
  \item 应最大化复用成熟发行资产。
\end{itemize}
\paragraph{Developer Workloads}
普通业务工作负载优先使用：

\begin{itemize}
  \item Plain YAML；
  \item Kustomize。
\end{itemize}
避免将 Helm Template Complexity 暴露给普通开发者。

\paragraph{Future Platform API}
长期演进方向为：

\begin{lstlisting}[language=yaml]
apiVersion: atlas.io/v1alpha1
kind: StreamingApplication
\end{lstlisting}
由平台 Operator 将高层领域模型展开成：

\begin{itemize}
  \item Flink；
  \item Kafka / Redpanda；
  \item Gateway；
  \item Observability；
  \item Policy；
\end{itemize}
等底层资源。

\vspace{0.5em}\hrule\vspace{0.5em}
\subsection{3. 标准目录与物理信任边界}
Atlas 的目录结构同时表达：

\begin{itemize}
  \item Domain Ownership；
  \item Trust Boundary；
  \item Agent Permission Boundary。
\end{itemize}
标准结构如下：

\begin{lstlisting}[language=text]
atlas/
├── env/                                      # Host Injection Context
│
├── bootstrap/                                # Phase 0 / Imperative Executor
│   └── argocd/
│       ├── install.sh
│       ├── atlas-bootstrap-project.yaml
│       └── root-app.yaml.tpl
│
├── gitops/                                   # Phase 1+ / Desired State
│   │
│   ├── root/                                 # 🔴 Tier-0 Trust Boundary
│   │   ├── base/
│   │   │   ├── kustomization.yaml
│   │   │   ├── project-bootstrap-app.yaml
│   │   │   ├── platform-control-app.yaml
│   │   │   └── workload-control-app.yaml
│   │   │
│   │   └── overlays/
│   │       ├── local-orbstack/
│   │       └── prod/
│   │
│   ├── environments/                         # Environment-specific K8s State
│   │   ├── local-orbstack/
│   │   └── prod/
│   │
│   ├── platform/                             # 🟡 Tier-1 Platform Baseline
│   │   │
│   │   ├── applications/                     # Platform Reconciliation Graph
│   │   │   ├── base/
│   │   │   └── overlays/
│   │   │       ├── local-orbstack/
│   │   │       └── prod/
│   │   │
│   │   ├── foundation/
│   │   │
│   │   ├── management/
│   │   │   ├── projects/
│   │   │   ├── argocd-self/
│   │   │   ├── sealed-secrets/
│   │   │   ├── cert-manager/
│   │   │   └── policy/
│   │   │
│   │   ├── operators/
│   │   ├── networking/
│   │   ├── observability/
│   │   ├── messaging/
│   │   └── storage/
│   │
│   └── workloads/                            # 🟢 Tier-2 Developer Domain
│       ├── applications/                     # Workload orchestration metadata
│       └── apps/                             # Actual workload desired state
│
├── vendor/                                   # Immutable Offline Artifacts
├── versions.lock                             # Shared Version Baseline
└── .state/                                   # Ephemeral Runtime State
    └── root-app.yaml
\end{lstlisting}
\vspace{0.5em}\hrule\vspace{0.5em}
\subsubsection{3.1 \_\_INLINE\_CODE\_0\_\_}
\_\_INLINE\_CODE\_0\_\_ 属于 Tier-0 Architecture Trust Boundary。

AI Agent 可以：

\begin{itemize}
  \item READ；
  \item ANALYZE；
  \item PROPOSE；
  \item GENERATE。
\end{itemize}
AI Agent MUST NOT 在无人类判断的情况下：

\begin{itemize}
  \item MERGE；
  \item APPLY；
  \item 修改 Root Trust Chain；
  \item 修改 Tier-0 AppProject Capability。
\end{itemize}
Tier-0 的自主修改必须经过 Human Judgment Gate。

\vspace{0.5em}\hrule\vspace{0.5em}
\subsubsection{3.2 \_\_INLINE\_CODE\_0\_\_}
该目录不是新的 Platform Domain。

它只表达：

\begin{gfmBlockquote}
Platform Application Graph / Reconciliation Metadata.
\end{gfmBlockquote}
其中只应保存：

\begin{itemize}
  \item Argo CD Application；
  \item Kustomization；
  \item Environment Overlay；
  \item 与 DAG 本身直接相关的调和元数据。
\end{itemize}
实际组件 Desired State MUST 位于其领域目录。

例如：

\begin{lstlisting}[language=text]
platform/applications/base/envoy-gateway-app.yaml
\end{lstlisting}
只定义 Application。

实际组件资产仍位于：

\begin{lstlisting}[language=text]
platform/networking/envoy-gateway/
\end{lstlisting}
因此：

\begin{lstlisting}[language=text]
WHERE does it belong?
        → Domain Directory

WHEN may it converge?
        → Application DAG
\end{lstlisting}
\vspace{0.5em}\hrule\vspace{0.5em}
\subsection{4. Atlas 控制拓扑}
Atlas 的逻辑控制图如下：

\begin{lstlisting}[language=text]
                    Git
                     │
                     │ Definition
                     ▼
                Bootstrap
                     │
                     │ Instantiation
                     ▼
          External Root Anchor
                     │
            Root Macro DAG
                     │
       ┌─────────────┼──────────────┐
       ▼             ▼              ▼
 Project Bootstrap  Platform      Workload
                    Control       Control
                       │              │
                       ▼              ▼
             Platform Capability   Workload
                    DAG              DAG
                       │              │
                       ▼              ▼
               Leaf Applications  Leaf Applications
                       │              │
                       └──────┬───────┘
                              ▼
                       Kubernetes API
                              ▲
                              │
                    Controllers / Operators
\end{lstlisting}
Root 与 Platform DAG 属于不同的同步作用域。

Sync Wave 不构成全平台共享的“绝对时间轴”。

其具体编号和 Health Gate 语义由：

\begin{gfmBlockquote}
Atlas GitOps Control Plane Design v1.0.3
\end{gfmBlockquote}
规范。

\vspace{0.5em}\hrule\vspace{0.5em}
\subsection{5. 工程约束与防御机制}
\vspace{0.5em}\hrule\vspace{0.5em}
\subsubsection{5.1 Dual Deletion Protection}
Atlas 将两类完全不同的灾难模型分开处理。

\paragraph{Git-side Accidental Deletion}
针对由 Parent Application 管理的 Tier-0 / Tier-1 Child Application CR：

\begin{lstlisting}[language=text]
argocd.argoproj.io/sync-options: Prune=confirm
\end{lstlisting}
Git 中误删除 Child Application 后：

\begin{lstlisting}[language=text]
Git deletion
    ↓
Parent detects prune
    ↓
Human confirmation required
\end{lstlisting}
\vspace{0.5em}\hrule\vspace{0.5em}
\paragraph{Explicit Application Deletion}
Atlas 的核心 Tier-0 / Tier-1 Application 默认：

\begin{gfmBlockquote}
MUST NOT carry cascading resources-finalizer.
\end{gfmBlockquote}
因此：

\begin{lstlisting}[language=text]
delete Application CR
       ↓
Application disappears
       ↓
Managed runtime resources survive
\end{lstlisting}
控制面事故不应自动演变为数据面销毁事故。

\vspace{0.5em}\hrule\vspace{0.5em}
\paragraph{5.1.1 Protection Scope}
Dual Deletion Protection 主要保护的是：

\begin{gfmBlockquote}
\textbf{Application Control Graph}
\end{gfmBlockquote}
它并不自动保护所有 descendant resources。

例如：

\begin{lstlisting}[language=text]
Redpanda Application
        ↓
StatefulSet / PVC / Cluster CR
\end{lstlisting}
如果 Leaf Application 对其资源执行 prune，Parent Application 上针对 Child Application CR 的 \_\_INLINE\_CODE\_0\_\_ 不会自动保护其内部 StatefulSet。

高价值资源必须根据领域需求额外采用：

\begin{itemize}
  \item resource-level \_\_INLINE\_CODE\_0\_\_；
  \item \_\_INLINE\_CODE\_0\_\_；
  \item Admission Policy；
  \item Domain-native protection；
  \item Backup / Restore。
\end{itemize}
因此：

\begin{gfmBlockquote}
Control Graph Deletion Protection ≠ Managed Resource Deletion Protection.
\end{gfmBlockquote}
\vspace{0.5em}\hrule\vspace{0.5em}
\subsubsection{5.2 Bootstrap Adoption Invariant}
Bootstrap Seed 与 \_\_INLINE\_CODE\_0\_\_ MUST 满足无损接管契约：

\begin{lstlisting}[language=text]
Object Identity Parity
+
Version Parity
+
Health Capability Parity
+
CRD Compatibility
=
Bootstrap Adoption Contract
\end{lstlisting}
\paragraph{Object Identity Parity}
Namespace、核心对象身份和关键资源命名必须兼容。

\paragraph{Version Parity}
Bootstrap 与 \_\_INLINE\_CODE\_0\_\_ MUST 使用同一份：

\begin{lstlisting}[language=text]
versions.lock
\end{lstlisting}
作为版本权威。

\paragraph{Health Capability Parity}
Seed 与 Full Desired State MUST 注入相同的 Application Health Capability。

\paragraph{CRD Compatibility}
恢复 Seed 不得造成：

\begin{itemize}
  \item CRD downgrade；
  \item Schema incompatibility；
  \item Controller binary / CRD mismatch。
\end{itemize}
Bootstrap 的目标不是：

\begin{gfmBlockquote}
安装一个“能跑”的 Argo CD。
\end{gfmBlockquote}
而是：

\begin{gfmBlockquote}
实例化一个能够被 Git Desired State 无损接管的最小控制面。
\end{gfmBlockquote}
\vspace{0.5em}\hrule\vspace{0.5em}
\subsubsection{5.3 Zero-Plaintext Secret}
任何未加密的 Secret，包括 Base64 编码内容：

\begin{gfmBlockquote}
MUST NOT enter Git.
\end{gfmBlockquote}
\paragraph{Phase 1}
采用 Sealed Secrets：

\begin{lstlisting}[language=text]
kubeseal
   ↓
SealedSecret
   ↓
Git
\end{lstlisting}
Controller Private Key 属于 Trust Root，必须物理隔离备份。

同时产生如下强依赖：

\begin{lstlisting}[language=text]
Sealed Secrets Controller
          ↓ Healthy
SealedSecret Materialization
          ↓
Secret Consumers
\end{lstlisting}
任何 Secret Consumer 不得先于 Secret Materialization Capability Ready。

\paragraph{Phase 2A}
演进至：

\begin{lstlisting}[language=text]
External Secrets Operator
+
External Secret Manager
\end{lstlisting}
例如：

\begin{itemize}
  \item Vault；
  \item AWS Secrets Manager；
  \item GCP Secret Manager；
  \item 等价外部密钥系统。
\end{itemize}
\vspace{0.5em}\hrule\vspace{0.5em}
\subsubsection{5.4 Tier-0 Human Judgment}
Atlas 不允许 AI Agent 在平台最核心的信任根拥有完全自治权。

以下操作 MUST 经过 Human Judgment：

\begin{itemize}
  \item Root Trust Chain 修改；
  \item Tier-0 AppProject 权限扩张；
  \item Break-Glass Recovery；
  \item Destructive Data Migration；
  \item Secret Trust Root 操作；
  \item 其他能够跨越信任域的高爆炸半径操作。
\end{itemize}
AI-Native 并不意味着：

\begin{gfmBlockquote}
Agent controls everything.
\end{gfmBlockquote}
Atlas 的目标是：

\begin{gfmBlockquote}
Agent has maximum autonomy inside explicitly bounded safe domains.
\end{gfmBlockquote}
\vspace{0.5em}\hrule\vspace{0.5em}
\subsection{6. Unified Telemetry \& Observability}
Atlas 将统一观测能力视为平台基础能力，而非应用上线后的附加功能。

\vspace{0.5em}\hrule\vspace{0.5em}
\subsubsection{6.1 Metrics}
采用：

\begin{itemize}
  \item Kube-Prometheus-Stack；
  \item Prometheus；
  \item OpenMetrics。
\end{itemize}
覆盖：

\begin{itemize}
  \item Kubernetes Control Plane；
  \item Platform Components；
  \item Operators；
  \item Application Services。
\end{itemize}
\vspace{0.5em}\hrule\vspace{0.5em}
\subsubsection{6.2 Collection}
统一采集代理采用：

\begin{itemize}
  \item Grafana Alloy；
  \item OpenTelemetry Collector。
\end{itemize}
\vspace{0.5em}\hrule\vspace{0.5em}
\subsubsection{6.3 Logs}
采用 Loki 作为日志后端，解决：

\begin{itemize}
  \item Pod 重建；
  \item Node 漂移；
  \item Container 生命周期结束；
\end{itemize}
后的历史追溯问题。

\vspace{0.5em}\hrule\vspace{0.5em}
\subsubsection{6.4 Traces}
采用 Tempo 作为 Trace Backend。

优先使用：

\begin{lstlisting}[language=text]
OTLP
\end{lstlisting}
作为跨服务遥测传输协议。

包括：

\begin{itemize}
  \item Envoy Gateway；
  \item Platform Services；
  \item Streaming Workloads；
  \item Microservices。
\end{itemize}
\vspace{0.5em}\hrule\vspace{0.5em}
\subsection{7. Disaster Recovery Boundaries}
Atlas 严格区分：

\begin{gfmBlockquote}
Infrastructure Recreation
与
Stateful Business Recovery
\end{gfmBlockquote}
\vspace{0.5em}\hrule\vspace{0.5em}
\subsubsection{7.1 Disposable State}
以下状态默认可丢弃：

\begin{itemize}
  \item Stateless Pod；
  \item Controller Runtime Cache；
  \item Git 可重建的 Kubernetes Objects；
  \item 临时执行环境；
  \item Bootstrap \_\_INLINE\_CODE\_0\_\_。
\end{itemize}
这些对象应通过 GitOps 或 Controller 快速重建。

\vspace{0.5em}\hrule\vspace{0.5em}
\subsubsection{7.2 Must-Backup State}
以下资产必须拥有独立备份方案：

\paragraph{Git Repository}
Git 是 Desired State 唯一权威。

\paragraph{Cluster Trust Roots}
包括：

\begin{itemize}
  \item Sealed Secrets Private Key；
  \item TLS Root；
  \item 外部密钥系统恢复凭证；
  \item 其他平台 Trust Material。
\end{itemize}
\paragraph{Business Data}
包括：

\begin{itemize}
  \item Kafka / Redpanda Logs；
  \item Database Physical Data；
  \item MinIO Object Data；
  \item Stateful Application Data。
\end{itemize}
\vspace{0.5em}\hrule\vspace{0.5em}
\subsubsection{7.3 Application-consistent Recovery}
Atlas 明确规定：

\begin{gfmBlockquote}
Storage-level backup MUST NOT be treated as application-consistent recovery.
\end{gfmBlockquote}
例如：

\paragraph{Flink}
依赖：

\begin{itemize}
  \item Checkpoint；
  \item Savepoint。
\end{itemize}
\paragraph{Kafka / Redpanda}
依赖：

\begin{itemize}
  \item Broker Replication；
  \item Cluster Metadata；
  \item Operator-aware Recovery。
\end{itemize}
\paragraph{Database}
依赖：

\begin{itemize}
  \item Database-native Backup / Restore；
  \item WAL / Binlog / Snapshot Semantics。
\end{itemize}
\paragraph{Object Storage}
依赖：

\begin{itemize}
  \item Object-level consistency；
  \item Replication；
  \item Versioning；
  \item Domain-native recovery mechanism。
\end{itemize}
PVC Snapshot 只能作为底层能力，不得替代领域恢复语义。

\vspace{0.5em}\hrule\vspace{0.5em}
\subsection{8. Break-Glass Architecture Principle}
Atlas 在灾难情况下执行：

\begin{gfmBlockquote}
Freeze Outside-In, Recover Inside-Out.
\end{gfmBlockquote}
\subsubsection{Freeze}
\begin{lstlisting}[language=text]
External Root
      ↓ freeze
Child Control Applications
      ↓ freeze
Affected Reconciliation Graph
\end{lstlisting}
优先阻断错误传播。

\subsubsection{Recover}
\begin{lstlisting}[language=text]
Seed
 ↓
argocd-self
 ↓
Platform Capability DAG
 ↓
Workload Control
 ↓
External Root Resume
\end{lstlisting}
优先恢复最内层可信控制能力，再逐级恢复外层调和。

具体 Break-Glass Procedure 由 GitOps Control Plane Design 管辖。

\vspace{0.5em}\hrule\vspace{0.5em}
\subsection{9. Architecture Invariants}
Atlas v1.0.2 正式冻结以下平台级不变量。

\subsubsection{INV-01 — Definition Authority}
Git MUST remain the authoritative definition of Desired State.

\subsubsection{INV-02 — Bootstrap Termination}
Bootstrap MUST cease active control after GitOps adoption.

\subsubsection{INV-03 — External Root}
External Root Anchor MUST remain outside parent Application reconciliation.

\subsubsection{INV-04 — GitOps Authority}
Argo CD MUST remain the sole continuous GitOps Desired-State reconciler.

\subsubsection{INV-05 — Runtime Authority}
Runtime lifecycle MUST belong to Kubernetes Controllers or Domain Operators.

\subsubsection{INV-06 — Two-Level DAG}
Atlas MUST use a minimal Root Macro DAG plus delegated capability-level DAGs.

\subsubsection{INV-07 — Root Minimalism}
Ordinary platform component lifecycle changes MUST NOT require Tier-0 Root changes.

\subsubsection{INV-08 — Domain / Dependency Separation}
Directory structure MUST represent domain ownership; dependency ordering MUST be represented separately.

\subsubsection{INV-09 — Leaf Failure Isolation}
Independent platform capabilities SHOULD map to independent Leaf Applications.

\subsubsection{INV-10 — Bounded Application Depth}
Additional App-of-Apps nesting MUST NOT be introduced solely to mirror directory hierarchy.

\subsubsection{INV-11 — Adoption Parity}
Bootstrap Seed and argocd-self MUST satisfy the Bootstrap Adoption Contract.

\subsubsection{INV-12 — Control/Data Failure Isolation}
Deleting or damaging the GitOps control graph MUST NOT implicitly destroy high-value runtime/data-plane state.

\subsubsection{INV-13 — Zero Plaintext Secret}
Unencrypted Secret material MUST NOT enter Git.

\subsubsection{INV-14 — Offline Artifact Availability}
All runtime-critical external artifacts MUST be locally available before entering offline operation.

\subsubsection{INV-15 — Stateful Recovery Semantics}
Application-consistent recovery MUST be owned by the corresponding stateful domain.

\subsubsection{INV-16 — Tier-0 Human Judgment}
Trust-boundary-changing operations MUST remain human-gated unless explicitly authorized by a future governance standard.

\vspace{0.5em}\hrule\vspace{0.5em}
\subsection{10. Conformance Model}
任何 Atlas 实现只有在同时满足：

\begin{lstlisting}[language=text]
Architecture Invariants
        +
GitOps Domain Contracts
        +
Repository Structural Rules
        +
Runtime Policy Controls
\end{lstlisting}
时，才可以声明：

\begin{gfmBlockquote}
Atlas Architecture Conformant.
\end{gfmBlockquote}
后续仓库审计应以：

\begin{lstlisting}[language=text]
INV-*
MUST
MUST NOT
SHOULD
\end{lstlisting}
为机械审查基准，而不是依赖自然语言主观解释。

\vspace{0.5em}\hrule\vspace{0.5em}
\subsection{11. Architecture Freeze Statement}
Atlas Architecture Design v1.0.2 正式冻结以下核心架构：

\begin{gfmBlockquote}
\textbf{External Root Anchor + Two-Level Reconciliation DAG}
\end{gfmBlockquote}
其中：

\begin{itemize}
  \item Root 只管理 Trust Transition；
  \item Platform Control 管理 Capability Dependency；
  \item Workload Control 管理 Workload Orchestration；
  \item Leaf Application 构成最小 Failure Isolation Unit；
  \item Domain Directory 与 Dependency Wave 完全解耦；
  \item Bootstrap、Argo CD、Operators 保持严格控制权边界。
\end{itemize}
未来普通平台组件的增加、删除与升级：

\begin{gfmBlockquote}
SHOULD NOT require modification of the Tier-0 Root Architecture.
\end{gfmBlockquote}
任何破坏上述模式的变更均视为 Architecture Change，必须通过正式 ADR / Architecture Review，而不得作为普通实现重构直接合入。

\end{document}